\documentclass{article}

\usepackage{neurips_2025}

\usepackage[utf8]{inputenc}
\usepackage[T1]{fontenc}
\usepackage{hyperref}
\usepackage{url}
\usepackage{booktabs}
\usepackage{amsfonts}
\usepackage{nicefrac}
\usepackage{microtype}
\usepackage{xcolor}
\usepackage{graphicx}

\title{Do Chess Explanations Encode Decisions?\\
Evaluating Human and LLM Reasoning with Move Recoverability}

\author{
Anonymous Author(s)
}

\begin{document}

\maketitle

\begin{abstract}
Large language models (LLMs) can generate fluent explanations for chess moves, but it remains unclear whether these explanations reflect the actual decision process or merely provide post-hoc narratives. 
We propose a framework for evaluating reasoning quality using \textit{move recoverability}: the ability to infer the intended move from an explanation. 
We conduct a pilot study using pawn endgame puzzles and human instructional transcripts from chess videos. 
Across six puzzles, LLM explanations show low recoverability even when prompted for step-by-step reasoning. 
In contrast, human instructional explanations contain structured reasoning such as conditional variations and named endgame concepts. 
Our findings suggest that LLM explanations frequently lack the causal structure necessary to reconstruct the underlying decision.
\end{abstract}

\section{Introduction}

Large language models increasingly generate explanations for decisions across many domains, including mathematics, programming, and strategy games. 
However, it remains unclear whether these explanations reflect genuine reasoning processes or simply produce plausible narratives after selecting an answer.

Chess provides a controlled testbed for studying explanation quality because:

\begin{itemize}
\item The correct move is objectively defined.
\item Move quality can be evaluated with chess engines.
\item Human reasoning can be observed through instructional commentary and think-aloud transcripts.
\end{itemize}

In this paper we investigate whether explanations encode sufficient information to reconstruct the intended move.

We introduce the metric \textbf{move recoverability}: given a position and an explanation (with explicit move references removed), can a decoder infer the intended move?

We evaluate recoverability for explanations generated by LLMs under different prompting strategies and compare them with human reasoning transcripts from instructional chess videos.

\section{Method}

\subsection{Puzzle Dataset}

We digitized puzzles from the endgame training book \textit{Test Your Endgame Ability}. 
For this pilot study we used six pawn endgame puzzles from the first section.

Each puzzle includes:

\begin{itemize}
\item FEN representation of the position
\item book solution move
\item Stockfish evaluation
\end{itemize}

\subsection{LLM Generation}

For each puzzle we generated explanations using three prompting styles:

\begin{enumerate}
\item \textbf{Brief explanation} (standard explanation)
\item \textbf{Calculation reasoning} (step-by-step analysis)
\item \textbf{Teaching explanation} (pedagogical explanation)
\end{enumerate}

For each generation we recorded:

\begin{itemize}
\item predicted move
\item legality
\item Stockfish evaluation difference
\item explanation text
\end{itemize}

\subsection{Move Recoverability}

To evaluate whether explanations encode the move decision we introduce the metric:

\textbf{Move Recoverability}

\begin{quote}
Given a chess position and a masked explanation, can a decoder infer the intended move?
\end{quote}

Explicit move mentions were removed from explanations before evaluation.

Recoverability is computed as:

\[
Recoverability = \frac{\text{correctly inferred moves}}{\text{total explanations}}
\]

\subsection{Human Reasoning Data}

To compare LLM explanations with human reasoning, we collected transcripts from instructional chess videos discussing pawn endgames.

Human explanations often include:

\begin{itemize}
\item conditional reasoning ("if the king moves here...")
\item named concepts (e.g., opposition, distant opposition)
\item explicit variation lines
\end{itemize}

Two example transcripts were manually annotated for this pilot study.

\section{Preliminary Results}

\subsection{LLM Puzzle Accuracy}

Table \ref{tab:accuracy} summarizes LLM performance on the six puzzles.

\begin{table}[h]
\centering
\begin{tabular}{lc}
\toprule
Metric & Value \\
\midrule
Legal move rate & 1.00 \\
Exact match rate & 0.33 \\
Average engine delta & 235 cp \\
\bottomrule
\end{tabular}
\caption{LLM performance on six pawn endgame puzzles.}
\label{tab:accuracy}
\end{table}

\subsection{Recoverability}

Recoverability results for each prompt style are shown in Table \ref{tab:recoverability}.

\begin{table}[h]
\centering
\begin{tabular}{lc}
\toprule
Prompt Style & Recoverability \\
\midrule
Brief & 0.33 \\
Calculation & 0.17 \\
Teaching & 0.00 \\
\bottomrule
\end{tabular}
\caption{Move recoverability across prompt styles.}
\label{tab:recoverability}
\end{table}

Even when prompted to produce step-by-step reasoning, explanations frequently lacked enough information to reconstruct the intended move.

\subsection{Qualitative Analysis}

Inspection of explanations reveals recurring narrative templates such as:

\begin{quote}
``advance the pawn toward promotion'' \\
``support the pawn'' \\
``restrict the opponent's king''
\end{quote}

These descriptions often capture general strategic ideas but fail to uniquely identify the correct move.

Human instructional explanations, by contrast, frequently contain structured reasoning patterns:

\begin{itemize}
\item conditional analysis of alternative moves
\item references to known endgame concepts
\item explicit variation trees
\end{itemize}

For example, human instructors frequently explain pawn endgames using the concept of \textit{opposition}, describing how the king must move to maintain control of key squares.

\section{Discussion}

The pilot results suggest that LLM explanations often function as post-hoc narratives rather than faithful reasoning traces.

Interestingly, prompting models to produce step-by-step calculations reduced recoverability relative to brief explanations.

One possible explanation is that LLMs generate plausible calculation trees that are not grounded in the actual board state.

\section{Future Work}

Future work will expand the dataset in several directions:

\begin{itemize}
\item scaling puzzle experiments to 100+ endgame positions
\item collecting spontaneous human think-aloud reasoning
\item evaluating human and LLM performance as explanation decoders
\item analyzing reasoning structure (conditional density, square mentions, concept vocabulary)
\end{itemize}

Chess puzzles provide a promising environment for studying reasoning transparency because both the decision outcome and evaluation function are well defined.

\section{Conclusion}

We introduce move recoverability as a metric for evaluating whether explanations encode the decision process behind chess moves.

In a small pilot study, LLM explanations exhibit low recoverability compared to structured human reasoning transcripts.

These findings suggest that fluent explanations do not necessarily reflect the mechanisms that produced the decision.

\section*{References}

{\small

[1] August Livshitz and Jon Speelman. \textit{Test Your Endgame Ability}. Chess Endgame Training.

[2] Herbert A. Simon and William Chase. Perception in chess. \textit{Cognitive Psychology}, 1973.

[3] Noam Brown et al. Superhuman AI for multiplayer poker. \textit{Science}, 2019.

}

\end{document}